%Version 3.1 December 2024
% See section 11 of the User Manual for version history
%
%%%%%%%%%%%%%%%%%%%%%%%%%%%%%%%%%%%%%%%%%%%%%%%%%%%%%%%%%%%%%%%%%%%%%%
%%                                                                 %%
%% Please do not use \input{...} to include other tex files.       %%
%% Submit your LaTeX manuscript as one .tex document.              %%
%%                                                                 %%
%% All additional figures and files should be attached             %%
%% separately and not embedded in the \TeX\ document itself.       %%
%%                                                                 %%
%%%%%%%%%%%%%%%%%%%%%%%%%%%%%%%%%%%%%%%%%%%%%%%%%%%%%%%%%%%%%%%%%%%%%

%%\documentclass[referee,sn-basic]{sn-jnl}% referee option is meant for double line spacing

%%=======================================================%%
%% to print line numbers in the margin use lineno option %%
%%=======================================================%%

%%\documentclass[lineno,pdflatex,sn-basic]{sn-jnl}% Basic Springer Nature Reference Style/Chemistry Reference Style

%%=========================================================================================%%
%% the documentclass is set to pdflatex as default. You can delete it if not appropriate.  %%
%%=========================================================================================%%

%%\documentclass[sn-basic]{sn-jnl}% Basic Springer Nature Reference Style/Chemistry Reference Style

%%Note: the following reference styles support Namedate and Numbered referencing. By default the style follows the most common style. To switch between the options you can add or remove “Numbered in the optional parenthesis. 
%%The option is available for: sn-basic.bst, sn-chicago.bst%  
 
%%\documentclass[pdflatex,sn-nature]{sn-jnl}% Style for submissions to Nature Portfolio journals
%%\documentclass[pdflatex,sn-basic]{sn-jnl}% Basic Springer Nature Reference Style/Chemistry Reference Style
\documentclass[lineno,pdflatex,sn-mathphys-num]{sn-jnl}% Math and Physical Sciences Numbered Reference Style
%%\documentclass[pdflatex,sn-mathphys-ay]{sn-jnl}% Math and Physical Sciences Author Year Reference Style
%%\documentclass[pdflatex,sn-aps]{sn-jnl}% American Physical Society (APS) Reference Style
%%\documentclass[pdflatex,sn-vancouver-num]{sn-jnl}% Vancouver Numbered Reference Style
%%\documentclass[pdflatex,sn-vancouver-ay]{sn-jnl}% Vancouver Author Year Reference Style
%%\documentclass[pdflatex,sn-apa]{sn-jnl}% APA Reference Style
%%\documentclass[pdflatex,sn-chicago]{sn-jnl}% Chicago-based Humanities Reference Style

%%%% Standard Packages
%%<additional latex packages if required can be included here>

\usepackage{graphicx}%
\usepackage{multirow}%
\usepackage{amsmath,amssymb,amsfonts}%
\usepackage{amsthm}%
\usepackage{mathrsfs}%
\usepackage[title]{appendix}%
\usepackage{xcolor}%
\usepackage{textcomp}%
\usepackage{manyfoot}%
\usepackage{booktabs}%
\usepackage{algorithm}%
\usepackage{algorithmicx}%
\usepackage{algpseudocode}%
\usepackage{listings}%
\usepackage{xurl}
\usepackage{placeins}
%%%%

%%%%%=============================================================================%%%%
%%%%  Remarks: This template is provided to aid authors with the preparation
%%%%  of original research articles intended for submission to journals published 
%%%%  by Springer Nature. The guidance has been prepared in partnership with 
%%%%  production teams to conform to Springer Nature technical requirements. 
%%%%  Editorial and presentation requirements differ among journal portfolios and 
%%%%  research disciplines. You may find sections in this template are irrelevant 
%%%%  to your work and are empowered to omit any such section if allowed by the 
%%%%  journal you intend to submit to. The submission guidelines and policies 
%%%%  of the journal take precedence. A detailed User Manual is available in the 
%%%%  template package for technical guidance.
%%%%%=============================================================================%%%%

%% as per the requirement new theorem styles can be included as shown below
\theoremstyle{thmstyleone}%
\newtheorem{theorem}{Theorem}%  meant for continuous numbers
%%\newtheorem{theorem}{Theorem}[section]% meant for sectionwise numbers
%% optional argument [theorem] produces theorem numbering sequence instead of independent numbers for Proposition
\newtheorem{proposition}[theorem]{Proposition}% 
%%\newtheorem{proposition}{Proposition}% to get separate numbers for theorem and proposition etc.

\theoremstyle{thmstyletwo}%
\newtheorem{example}{Example}%
\newtheorem{remark}{Remark}%

\theoremstyle{thmstylethree}%
\newtheorem{definition}{Definition}%

\raggedbottom
%%\unnumbered% uncomment this for unnumbered level heads

% \usepackage{amsmath}    % for \text{} inside equations
% \usepackage{multirow}   % optional, for merged table cells
% \usepackage{tikz}
% \usepackage{url}
% \usetikzlibrary{arrows.meta,positioning}
% \usetikzlibrary{calc}
% \usetikzlibrary{patterns}
% \usepackage{pgfplots}   % bar charts


\begin{document}

\title[Article Title]{Where Do Vision Foundation Models Look? A Study of Attention Alignment in Image Recognition}

%%=============================================================%%
%% GivenName	-> \fnm{Joergen W.}
%% Particle	-> \spfx{van der} -> surname prefix
%% FamilyName	-> \sur{Ploeg}
%% Suffix	-> \sfx{IV}
%% \author*[1,2]{\fnm{Joergen W.} \spfx{van der} \sur{Ploeg} 
%%  \sfx{IV}}\email{iauthor@gmail.com}
%%=============================================================%%

\author[1]{\fnm{Desmond} \sur{Choy}}\email{e1510259@u.nus.edu}
\author[1]{\fnm{Kay Mei} \sur{Leong}}\email{e0324286@u.nus.edu}
\author*[1]{\fnm{Jing} \sur{Tian}}\email{tianjing@nus.edu.sg}
\author[1]{\fnm{Xavier Zhen} \sur{Xie}}\email{xav.xie@nus.edu.sg}

\affil*[1]{\orgdiv{NUS-ISS}, \orgname{National University of Singapore}, \orgaddress{\street{25 Heng Mui Keng Terrace}, \city{Singapore}, \postcode{119615}, \country{Singapore}}}


%%==================================%%
%% Sample for unstructured abstract %%
%%==================================%%

% \abstract{Vision foundation models achieve strong recognition performance, but it remains unclear whether their attention aligns with the visual evidence that human experts consider discriminative. This paper presents a comparative study of attention alignment in image recognition, using the WikiChurches benchmark dataset, in which expert-annotated bounding boxes identify style-relevant architectural features. The study examines the attention alignment by formulating three research questions: whether frozen models exhibit expert-aligned attention, how attention alignment varies across intermediate layers and heads, and how downstream adaptation changes attention alignment. This paper presents a comprehensive comparison of six vision foundation models, spanning image-only self-supervised and vision-language models, and measures attention alignment using five complementary metrics, three types of attention maps, and three adaptation strategies. The findings show that attention alignment provides a useful diagnostic lens beyond classification accuracy, while remaining evidence of spatial plausibility rather than causal explanation.}


% \keywords{Explainable computer vision; vision foundation model; attention map; attentional alignment; image recognition}

%%\pacs[JEL Classification]{D8, H51}

%%\pacs[MSC Classification]{35A01, 65L10, 65L12, 65L20, 65L70}

\maketitle



\section*{Declarations}

\noindent\textbf{Competing interests}

The authors declare no financial or commercial conflicts of interest. \newline

\noindent\textbf{Data availability }

The data that supports the findings of this study is available on the GitHub repository \url{https://github.com/keimiii/xai\_viz}.\newline

\noindent\textbf{Code availability }

The code that supports the findings of this study is available on the GitHub repository \url{https://github.com/keimiii/xai\_viz}.\newline


\noindent\textbf{Funding}

This research received no external funding. \newline


\noindent\textbf{Ethics approval and consent to participate}

Not applicable.\newline

\newpage 
\noindent\textbf{Consent for publication}

Not applicable.\newline



\noindent\textbf{Author contribution}

D. Choy: Conceptualization, Methodology, Software, Validation, Formal Analysis, Investigation, Data Curation, Writing - Original Draft, Writing - Review \& Editing, Visualization

K. Leong: Conceptualization, Methodology, Software, Validation, Formal Analysis, Writing - Original Draft, Writing - Review \& Editing

J. Tian: Formal analysis, Writing - Original Draft, Writing - Review \& Editing, Supervision.

X. Xie: Formal analysis, Writing - Original Draft, Writing - Review \& Editing.



\end{document}
